\documentclass{../base/canon}

%% canon_variables.tex — Valores por defecto de metadatos institucionales
%%
%% Incluir ANTES del \begin{document} para sobreescribir los defaults de la clase.
%% El template puede sobreescribir cualquier valor después de este \input.
%%
%% Uso en template:
%%   %% canon_variables.tex — Valores por defecto de metadatos institucionales
%%
%% Incluir ANTES del \begin{document} para sobreescribir los defaults de la clase.
%% El template puede sobreescribir cualquier valor después de este \input.
%%
%% Uso en template:
%%   %% canon_variables.tex — Valores por defecto de metadatos institucionales
%%
%% Incluir ANTES del \begin{document} para sobreescribir los defaults de la clase.
%% El template puede sobreescribir cualquier valor después de este \input.
%%
%% Uso en template:
%%   \input{../base/canon_variables}
%%   \SetDocTitle{Mi Informe Específico}   % sobreescribe el default

\SetDocLogo{../assets/logo.pdf}
\SetDocUnit{Tecnología \& Sistemas}
\SetContactUnit{Mesa de soporte}
\SetContactMembers{%
  Equipo CoreX & Soporte & soporte@empresa.com%
}
\SetContactNote{Para soporte técnico o consultas sobre este documento.}

%%   \SetDocTitle{Mi Informe Específico}   % sobreescribe el default

\SetDocLogo{../assets/logo.pdf}
\SetDocUnit{Tecnología \& Sistemas}
\SetContactUnit{Mesa de soporte}
\SetContactMembers{%
  Equipo CoreX & Soporte & soporte@empresa.com%
}
\SetContactNote{Para soporte técnico o consultas sobre este documento.}

%%   \SetDocTitle{Mi Informe Específico}   % sobreescribe el default

\SetDocLogo{../assets/logo.pdf}
\SetDocUnit{Tecnología \& Sistemas}
\SetContactUnit{Mesa de soporte}
\SetContactMembers{%
  Equipo CoreX & Soporte & soporte@empresa.com%
}
\SetContactNote{Para soporte técnico o consultas sobre este documento.}

\SetDocType{ACTA DE REUNIÓN}
\SetDocTitle{Revisión Operativa Semanal — Semana 17-2026}
\SetDocCode{ACT-OPS-028}
\SetDocArea{Operaciones / Postcosecha}
\SetDocAuthor{Patricia Suárez}
\SetDocDate{2026-04-22}
\SetContactUnit{Coordinación de Operaciones}
\SetContactMembers{%
  Patricia Suárez  & Coordinadora     & p.suarez@empresa.com \\
  Carlos Mendoza   & Jefe Operaciones & c.mendoza@empresa.com%
}

\begin{document}

\begin{table}[H]
  \begin{tabular}{@{}ll@{}}
    \textbf{Fecha:}    & Martes 2026-04-22 \\
    \textbf{Hora:}     & 08:00 -- 09:15 \\
    \textbf{Lugar:}    & Sala de reuniones — Planta Postcosecha \\
    \textbf{Objetivo:} & Revisión de resultados semana 17 y planificación semana 18 \\
  \end{tabular}
\end{table}

\section{Asistentes}

\begin{AsistentesBlock}
  Carlos Mendoza   & Jefe de Operaciones      & c.mendoza@empresa.com    \\
  Patricia Suárez  & Coordinadora de Calidad  & p.suarez@empresa.com     \\
  Roberto Flores   & Jefe de Tecnología       & r.flores@empresa.com     \\
  Diego Arévalo    & Ingeniero de Datos       & d.arevalo@empresa.com    \\
  María Quiñonez   & Supervisora Turno Mañana & m.quinonez@empresa.com   \\
\end{AsistentesBlock}

\section{Temas tratados}

\begin{enumerate}
  \item \textbf{Resultados semana 17.}  Carlos Mendoza presentó el informe
        \texttt{INF-OPS-017}.  El rendimiento de 96.3\,\% fue evaluado como positivo.
        Se discutió el incremento puntual de merma del turno tarde del día 15, que
        quedó explicado por la llegada de lotes de mayor madurez.

  \item \textbf{Gap de datos BAL2.}  Diego Arévalo informó sobre el gap de datos
        entre las 02:15 y las 04:40 del 2026-04-17 (ver \texttt{REP-TEC-003}).
        Se acordó instalar UPS como medida preventiva.  Roberto Flores confirmará
        disponibilidad del equipo antes del 2026-04-25.

  \item \textbf{Planificación semana 18.}  Se revisó el plan de producción
        provisional.  Se estiman 855\,000 tallos.  María Quiñonez señaló que el
        turno mañana requiere un operario adicional en apertura para los días lunes
        y martes por vacaciones programadas.

  \item \textbf{Mantenimiento ERP del 2026-04-24.}  Roberto Flores confirmó el
        mantenimiento de 20:00 a 24:00.  Se instruyó a los supervisores mediante el
        memorando \texttt{MEM-GER-001} ya distribuido.
\end{enumerate}

\section{Acuerdos y compromisos}

\begin{AcuerdosList}
  \item \textbf{R. Flores — 2026-04-25:} Confirmar disponibilidad de UPS para el
        servidor de adquisición BAL2 e indicar fecha de instalación.

  \item \textbf{C. Mendoza — 2026-04-23:} Gestionar con Recursos Humanos la
        cobertura del operario adicional en apertura para semana 18.

  \item \textbf{D. Arévalo — 2026-04-30:} Documentar el método de imputación
        \texttt{imp\_avg3} en el manual de operaciones del sistema de balanzas.

  \item \textbf{P. Suárez — 2026-04-30:} Proponer ajuste del umbral de merma por
        «botón abierto» para variedades de tallo corto al comité de calidad.

  \item \textbf{Todos los supervisores — 2026-04-24 19:30:} Completar cierre de
        registros antes del mantenimiento según instrucción \texttt{MEM-GER-001}.
\end{AcuerdosList}

\ObservationBox[Próxima reunión]{
  Martes 2026-04-29, 08:00, Sala de reuniones Planta Postcosecha.
  Orden del día: resultados semana 18 + seguimiento de acuerdos esta acta.
}

\section{Cierre}

La reunión concluyó a las \textbf{09:15}.  Los participantes se comprometen a
cumplir los acuerdos en los plazos indicados y a reportar impedimentos a la
coordinación con al menos 48 horas de anticipación.

\vspace{10pt}
\noindent
\begin{minipage}{0.45\linewidth}
  \SignatureBlock{Carlos Mendoza}{Jefe de Operaciones Postcosecha}
\end{minipage}
\hfill
\begin{minipage}{0.45\linewidth}
  \SignatureBlock{Patricia Suárez}{Secretaria de Acta}
\end{minipage}

\ContactSignature

\end{document}
